% GENERATED FILE. Edit canonical lessons, then run npm run generate:books.
% canonical-source-hash: fnv1a64:efc0279130e76e57
% canonical-lessons: GE-C121-bohnen, GE-C121-gurke, GE-C121-erdbeere, GE-C121-kirsche, GE-C121-melone

\chapter{Beans, Cucumber, Strawberry, Cherry, Melon}
\label{ch:ge-beans-cucumber-strawberry-cherry-melon}
\hlchaptermodality{121}

\begin{chapteropening}
\textbf{By the end of this chapter:} \emph{I can say beans, a cucumber, a strawberry, a cherry and a melon.}

Complete the last lesson of chapter 121: I can say beans, a cucumber, a strawberry, a cherry and a melon.
\end{chapteropening}

\section[die Bohnen]{die Bohnen — beans}

\label{lesson:GE-C121-bohnen}



\begin{quote}
Before the new one: say the German for vegetables, then the German for a mushroom.

Beans come in the plural: say \emph{die Bohnen} with its plural \emph{die}, and then \emph{die Jahre}.
\end{quote}

\subsection*{You'll want to know: die Bohnen}
\textbf{die Bohnen} — \textquotedblleft{}beans\textquotedblright{}.

\begin{cousinweb}
\textbf{Bohnen} is the plural of \emph{Bohne} $\leftarrow$ Old High German \emph{bōna} $\leftarrow$ Proto-Germanic \emph{*baunō} — inherited, the same word as English \emph{bean}. (EtymDB 2.1)
\end{cousinweb}

\begin{grammarlens}[title={What you've built}]
One more word for this chapter.
\end{grammarlens}

\subsection*{Your turn}
\begin{itemize}
\raggedright
  \item \emph{Say it:} \emph{die Bohnen}
  \item \emph{Say it:} \emph{die Bohnen}, once more
  \item \emph{Say it:} say \emph{die Bohnen}
  \item \emph{Recall:} say the German for vegetables, then the German for a mushroom, then say \emph{die Bohnen} again
\end{itemize}

\begin{tcolorbox}[breakable,colback=teal!4,colframe=teal!35!black,title={Before you move on}]
What does die Bohnen mean? (\textquotedblleft{}Beans\textquotedblright{}.) Say it once more.
\end{tcolorbox}
\section[die Gurke]{die Gurke — a cucumber}

\label{lesson:GE-C121-gurke}



\begin{quote}
Before the new one: say the German for a mushroom, then the German for beans.
\end{quote}

\subsection*{You'll want to know: die Gurke}
\textbf{die Gurke} — \textquotedblleft{}a cucumber\textquotedblright{}.

\begin{cousinweb}
\textbf{Gurke} came in from a West Slavic word, like Polish \emph{ogórek}, which goes back to Medieval Greek — borrowed. EtymDB 2.1 has the German word but not its source.
\end{cousinweb}

\begin{grammarlens}[title={What you've built}]
One more word for this chapter.
\end{grammarlens}

\subsection*{Your turn}
\begin{itemize}
\raggedright
  \item \emph{Say it:} \emph{die Gurke}
  \item \emph{Say it:} \emph{die Gurke}, once more
  \item \emph{Say it:} say \emph{die Gurke}
  \item \emph{Recall:} say the German for a mushroom, then the German for beans, then say \emph{die Gurke} again
\end{itemize}

\begin{tcolorbox}[breakable,colback=teal!4,colframe=teal!35!black,title={Before you move on}]
What does die Gurke mean? (\textquotedblleft{}A cucumber\textquotedblright{}.) Say it once more.
\end{tcolorbox}
\section[die Erdbeere]{die Erdbeere — a strawberry}

\label{lesson:GE-C121-erdbeere}



\begin{quote}
Before the new one: say the German for beans, then the German for a cucumber.
\end{quote}

\subsection*{You'll want to know: die Erdbeere}
\textbf{die Erdbeere} — \textquotedblleft{}a strawberry\textquotedblright{}.

\begin{cousinweb}
\textbf{Erdbeere} is the word for \textquotedblleft{}earth\textquotedblright{}, Old High German \emph{erda}, English \emph{earth}, + the word for \textquotedblleft{}berry\textquotedblright{}, English \emph{berry} (EtymDB 2.1): the berry that grows on the ground.
\end{cousinweb}

\begin{grammarlens}[title={What you've built}]
One more word for this chapter.
\end{grammarlens}

\subsection*{Your turn}
\begin{itemize}
\raggedright
  \item \emph{Say it:} \emph{die Erdbeere}
  \item \emph{Say it:} \emph{die Erdbeere}, once more
  \item \emph{Say it:} say \emph{die Erdbeere}
  \item \emph{Recall:} say the German for beans, then the German for a cucumber, then say \emph{die Erdbeere} again
\end{itemize}

\begin{tcolorbox}[breakable,colback=teal!4,colframe=teal!35!black,title={Before you move on}]
What does die Erdbeere mean? (\textquotedblleft{}A strawberry\textquotedblright{}.) Say it once more.
\end{tcolorbox}
\section[die Kirsche]{die Kirsche — a cherry}

\label{lesson:GE-C121-kirsche}



\begin{quote}
Before the new one: say the German for a cucumber, then the German for a strawberry.
\end{quote}

\subsection*{You'll want to know: die Kirsche}
\textbf{die Kirsche} — \textquotedblleft{}a cherry\textquotedblright{}.

\begin{cousinweb}
\textbf{Kirsche} $\leftarrow$ Old High German \emph{chirsa} $\leftarrow$ Vulgar Latin \emph{ceresia}, \textquotedblleft{}cherry\textquotedblright{} (EtymDB 2.1) — borrowed, like English \emph{cherry}.
\end{cousinweb}

\begin{grammarlens}[title={What you've built}]
One more word for this chapter.
\end{grammarlens}

\subsection*{Your turn}
\begin{itemize}
\raggedright
  \item \emph{Say it:} \emph{die Kirsche}
  \item \emph{Say it:} \emph{die Kirsche}, once more
  \item \emph{Say it:} say \emph{die Kirsche}
  \item \emph{Recall:} say the German for a cucumber, then the German for a strawberry, then say \emph{die Kirsche} again
\end{itemize}

\begin{tcolorbox}[breakable,colback=teal!4,colframe=teal!35!black,title={Before you move on}]
What does die Kirsche mean? (\textquotedblleft{}A cherry\textquotedblright{}.) Say it once more.
\end{tcolorbox}
\section[die Melone]{die Melone — a melon}

\label{lesson:GE-C121-melone}



\begin{quote}
Before the new one: say the German for a strawberry, then the German for a cherry.
\end{quote}

\subsection*{You'll want to know: die Melone}
\textbf{die Melone} — \textquotedblleft{}a melon\textquotedblright{}.

\begin{cousinweb}
\textbf{Melone} $\leftarrow$ Italian \emph{melone} $\leftarrow$ Latin \emph{mēlo}, from Greek \emph{mēlopépōn}, \textquotedblleft{}apple-melon\textquotedblright{} — borrowed. EtymDB 2.1 has no row for the German word.
\end{cousinweb}

\begin{grammarlens}[title={What you've built}]
One more word for this chapter.
\end{grammarlens}

\subsection*{Your turn}
\begin{itemize}
\raggedright
  \item \emph{Say it:} \emph{die Melone}
  \item \emph{Say it:} \emph{die Melone}, once more
  \item \emph{Say it:} say \emph{die Melone}
  \item \emph{Recall:} say the German for a strawberry, then the German for a cherry, then say \emph{die Melone} again
\end{itemize}

\begin{tcolorbox}[breakable,colback=teal!4,colframe=teal!35!black,title={Before you move on}]
What does die Melone mean? (\textquotedblleft{}A melon\textquotedblright{}.) Say it once more.
\end{tcolorbox}
